\chapter{量子涨落与结构的起源}

\lectureoverview{最后一讲要问的是：一个被暴胀抹平了的宇宙，如何仍能获得星系赖以生长的微小不规则性。答案分阶段建立。普适常数首先定义了普朗克尺度，并揭示出熟悉的粒子物理与量子引力之间巨大的鸿沟。微波天空随后提供了目标：量级为\(10^{-5}\)的温度与密度变化，分布在许多尺度上。零点运动提供了不可避免的涨落，阻尼解释了振荡如何能够冻结，而膨胀宇宙中的波动方程则把暴胀子的每一个傅里叶模式都变成恰好这样的阻尼振子。当模式的物理波长达到哈勃尺度时，模式冻结。暴胀结束时，冻结的变化使相邻区域在略微不同的时刻再加热，从而把微观量子噪声转化为原初密度图案。}

\section{哪些常数是普适的？}

我们拥有一个关于基本粒子的极其精确的量子理论。它的参数包括夸克与轻子的质量、耦合常数、混合数据、电荷，以及粒子本身的清单。给定无限的计算能力，这个理论不仅能解释加速器实验，还能解释原子、分子和化学；原则上，其后果一直延伸到生物学以及思维背后的物理过程。这并不意味着这类系统容易计算。它说的是，标准模型意在作为普通物质的微观基础。引力则另行描述，在经典层面上由广义相对论给出。两个理论都很强大；但两者都没有使低能物理中出现的每一个数字显得不可避免。

质子与电子的质量比是一个完全良好的物理数字，但它与两个特定的粒子绑定。把质子换成中子，或把电子换成{\proseunicodefont μ}子，就会产生不同的比值。这类参数或许编码着我们尚不了解的深层物理，但它们并不彰显自身的普适性。三个量纲常数则不同：
\[
c,\qquad \hbar,\qquad G.
\]
它们的定义定律中反复出现\emph{一切}一词。一切物质与信息都遵守由\(c\)给出的极限速度。一切量子系统都服从由\(\hbar\)支配的不确定性关系。一切形式的能量都以同一个牛顿常数\(G\)产生引力。正是这种广度，使这些常数在当前的论证中成为基本的。

普朗克在量子力学的完整含义尚未为人所知之前，就认识到了这一组合的重要性。他的常数最初进入黑体辐射的描述。即便在那时，量纲上的事实也已显而易见：从\(c\)、\(\hbar\)和\(G\)出发，无需参照任何特定的原子或粒子，就能构造出长度、时间和质量的单位。

\section{构造普朗克尺度}

设\(L\)、\(T\)和\(M\)分别表示长度、时间和质量的量纲。光速的量纲为
\[
[c]=LT^{-1}.
\]
不确定性关系\(\Delta x\,\Delta p\gtrsim\hbar\)连同\(p=mv\)给出
\[
[\hbar]=ML^2T^{-1}.
\]
最后，牛顿方程
\[
F=\frac{Gm_1m_2}{r^2},
\qquad
F=ma
\]
给出
\[
[G]=L^3M^{-1}T^{-2}.
\]
这一推导之所以重要，是因为量纲分析并非猜测：每个指数都由要求不需要的量纲相消而确定。

寻求如下形式的长度
\begin{equation}
\ell=c^\alpha\hbar^\beta G^\gamma.
\label{eq:c10-length-ansatz}
\end{equation}
质量的指数必须为零，因此
\[
\beta-\gamma=0.
\]
时间的指数必须为零，而长度的指数必须等于一：
\[
-\alpha-\beta-2\gamma=0,
\qquad
\alpha+2\beta+3\gamma=1.
\]
利用\(\gamma=\beta\)，这些式子变为
\[
\alpha+3\beta=0,
\qquad
\alpha+5\beta=1.
\]
两者相减给出\(2\beta=1\)，因此
\[
\beta=\gamma=\frac12,
\qquad
\alpha=-\frac32.
\]
所以，由这三个普适常数构成的唯一长度为
\begin{equation}
\ell_{\mathrm P}
=\sqrt{\frac{\hbar G}{c^3}}
\approx 1.6\times10^{-35}\,\mathrm m.
\label{eq:c10-planck-length}
\end{equation}
本讲只保留其数量级\(10^{-35}\,\mathrm m\)，但量纲推导确定了这一组合本身。

与之对应的时间是光穿越时间：
\begin{equation}
t_{\mathrm P}
=\frac{\ell_{\mathrm P}}{c}
=\sqrt{\frac{\hbar G}{c^5}}
\approx 5.4\times10^{-44}\,\mathrm s.
\label{eq:c10-planck-time}
\end{equation}
在数量级的精度下，这就是\(10^{-43}\,\mathrm s\)。一旦普朗克长度和普适速度已知，它便不再是一个独立的构造。

对于质量，利用量纲关系\(Et\sim\hbar\)和\(E=mc^2\)。取\(t=t_{\mathrm P}\)给出
\begin{equation}
m_{\mathrm P}
=\frac{\hbar}{c^2t_{\mathrm P}}
=\sqrt{\frac{\hbar c}{G}}
\approx 2.2\times10^{-8}\,\mathrm{kg}.
\label{eq:c10-planck-mass}
\end{equation}
以粒子的标准来看，这大得惊人地宏观，大致相当于一粒细小尘埃的质量。其静能约为\(10^9\,\mathrm J\)，在量级上与一箱燃料的化学能相当。一辆汽车只需几箱燃料就能行驶数百英里，因为燃烧仅仅释放出燃料静能中极小的一部分。把一个普朗克质量完全转化为能量则完全是另一回事。当乘以\(c^2\)时，微小的质量可以代表巨大的能量。

与加速器物理的对比十分鲜明。在讲座当时的高能碰撞中所探测的最短尺度，其数量级为\(10^{-18}\,\mathrm m\)，这仍然大约是\(10^{17}\)个普朗克长度。因此，电弱尺度与普朗克尺度之间的未知领域极为广阔。核物理提供了恰当的类比。从远处看，原子核像是基本粒子，但在更短的距离上，它分解为质子和中子，而这些又进一步分解为夸克和胶子。低能粒子与原子核相比可能显得简单，但相对于普朗克尺度上存在的任何东西而言，它们仍然是复杂的涌现客体。从这一层级结构来看，标准模型的众多参数就不那么令人惊讶了，尽管其观测值和精细调节仍然没有得到解释。

\begin{classroomqa}
\audiencequestion{希格斯粒子比预期的更轻还是更重，为什么它的质量令人困惑？}
\lecturerresponse{有两种比较。相对于普朗克质量，一个\(125\,\mathrm{GeV}\)的希格斯粒子极其轻。在没有保护结构的理论中，量子修正往往使标量质量对高得多的尺度敏感，因此要得到观测值看起来就像在狭窄的边缘上平衡参数。相对于当时讨论的某些简单超对称预期，观测到的希格斯粒子反而可以说有些重。表面上的矛盾来自将其与两个不同的参考尺度进行比较；底层的层级问题依然存在。\lecturetimestamp{00:26:45--00:30:54}}
\end{classroomqa}

\begin{classroomqa}
\audiencequestion{如果希格斯粒子帮助赋予粒子质量，那么暗物质粒子的质量是否应该接近普朗克尺度？}
\lecturerresponse{不。WIMP是\emph{弱相互作用大质量粒子}；这个说法指向的是一种可能使实验室探测成为可能的相互作用尺度，而不是普朗克质量。在通常的WIMP假设下，质量可以远高于质子质量，同时仍然远低于\(m_{\mathrm P}\)。如果允许与普通物质的相互作用任意弱，宇宙学允许的质量范围会宽得多，但直接发现相应地变得更加困难。质子和中子也提供了一个有用的警示：它们的大部分质量来自强相互作用的束缚能，而不仅仅是希格斯机制产生的夸克质量之和。\lecturetimestamp{00:30:55--00:33:11}}
\end{classroomqa}

\begin{classroomqa}
\audiencequestion{普朗克质量意味着什么，是否存在每比特的基本能量？}
\lecturerresponse{一种有用的引力解释是，黑洞不可能在保持常规半经典黑洞的同时，在参数上被做得比普朗克尺度更小。这样的客体具有\(m_{\mathrm P}\)量级的质量、一个普朗克面积量级的视界面积，以及一个单位量级的熵。在这种受限的黑洞意义上，人们可以将普朗克量级的能量与一阶的信息量联系起来。这并不是一个普遍性论断，即每个设备中存储的每一比特都耗费一个普朗克能量；普通粒子和普通存储器可以轻得多。\lecturetimestamp{00:33:11--00:35:10}}
\end{classroomqa}

从这里开始，使用普朗克单位会很方便，
\begin{equation}
c=\hbar=G=1.
\label{eq:c10-planck-units}
\end{equation}
于是长度、时间、质量和能量都可以表示为相对于方程~\eqref{eq:c10-planck-length}--\eqref{eq:c10-planck-mass}的纯数。

\section{几乎均匀的微波天空}

宇宙学问题始于一张全天微波图。WMAP，即威尔金森微波各向异性探测器，测量了在原始等离子体变得透明的那一时期前后释放的辐射。这些光子此后被红移了大约一千倍，现在形成温度约为\(2.7\,\mathrm K\)的几乎完美的黑体。当我们向外看时，我们看到的并不是处于同一个当前时刻的宇宙。距离越远意味着发射越早，而微波背景就是光子首次自由地朝我们流来的遥远壳层。

\begin{classroomqa}
\audiencequestion{椭圆形的WMAP图像究竟代表什么？}
\lecturerresponse{椭圆形中的每一点代表天空中的一个方向，球面通过地图投影被展平。左右边缘相接，上下部分代表相反的银道半球，银道面横贯中间。来自我们自己银河系的前景微波辐射必须被建模并扣除。剩下的不是附近恒星的图，而是最后散射壳层上微波温度的方向图。\lecturetimestamp{00:37:49--00:40:49}}
\end{classroomqa}

在一级近似下，结果是各向同性的：
\[
T(\hat{\mathbf n})\simeq 2.7\,\mathrm K
\]
对每个观测方向\(\hat{\mathbf n}\)都成立。围绕每一点的完美各向同性将意味着完美均匀性。但一个完美均匀的宇宙不会自发产生星系、星系团以及上一讲讨论过的类似焦散的宇宙网。必定存在某种非均匀性，无论其幅度多么小。

将方向性温度偏离写为
\[
\delta T(\hat{\mathbf n})=T(\hat{\mathbf n})-\bar T.
\]
测得的相对变化量约为
\begin{equation}
\frac{\delta T}{T}\sim10^{-5}.
\label{eq:c10-cmb-anisotropy}
\end{equation}
得到这个数值的一条途径是动力学途径：将今天的物质分布反向演化，并追问在解耦时需要多大的初始对比度。另一条途径是直接的：将几乎无特征的微波图放大五个数量级，并测量其斑驳本身。这两条途径在本论证所需的精度上一致。

斑驳出现在许多角尺度上。大块区域包含较小的区域，而放大某一区域会得到统计上相似的纹理。\emph{尺度不变}一词并不意味着每个尺度都字面上完全相同；它意味着在相关范围内，涨落功率随尺度的变化很弱。这种近尺度不变性是理论必须解释的证据的一部分，而不仅仅是整体幅度。

为什么温度应当反映物质？一个光子逃离引力势阱会损失能量。因此，来自较密区域的辐射相对于来自较疏区域的辐射会发生偏移。在本讲的量级水平上，这为解耦附近
\begin{equation}
\frac{\delta\rho}{\rho}\sim10^{-5}
\label{eq:c10-density-target}
\end{equation}
提供了动机。精确的CMB计算包括声学动力学、引力势和转移函数，因此方程~\eqref{eq:c10-density-target}是设定尺度的简写，而不是\(\delta T/T\)与\(\delta\rho/\rho\)之间的精确等式。

\begin{classroomqa}
\audiencequestion{斑驳会不会是最后散射时产生的普通微观噪声？}
\lecturerresponse{在观测到的角尺度上不会。到解耦时，这些斑块对应着巨大的物理区域，而不是复合等离子体中原子尺度的扰动。然而，将它们追溯得更远得多，会使其物理波长变小。暴胀提供了缺失的放大机制：起初作为量子、微观扰动的涨落可以被拉伸到天文尺度。\lecturetimestamp{00:46:42--00:48:02}}
\end{classroomqa}

该机制可以从三部分物理来理解：谐振子的零点运动、振子的阻尼，以及膨胀宇宙中标量场的波动方程。我们将按此顺序把它们组合起来。

\section{振子的零点运动}

对于单位质量的谐振子，
\begin{equation}
E=\frac12\dot x^{\,2}+\frac12\omega^2x^2,
\qquad
\ddot x+\omega^2x=0.
\label{eq:c10-undamped}
\end{equation}
动能和势能不断交换。在一次振荡上取平均，二者各占总能量的一半。

量子力学给出能级
\begin{equation}
E_n=\left(n+\frac12\right)\hbar\omega.
\label{eq:c10-oscillator-levels}
\end{equation}
基态不能同时具有\(x=0\)和\(\dot x=0\)，因为那将同时指定精确的位置和动量。其不可约能量
\[
E_0=\frac12\hbar\omega,
\]
就是零点能。\emph{振荡}一词应从统计上理解：定态基态并不描绘经典的 sinusoidal 轨道，但它在位置和动量上具有非零方差。

平均而言，基态能量的一半在势能项中：
\[
\frac12\omega^2\langle x^2\rangle
=\frac14\hbar\omega.
\]
因此
\begin{equation}
\langle x^2\rangle=\frac{\hbar}{2\omega}
\sim\frac{\hbar}{\omega}.
\label{eq:c10-zero-point-position}
\end{equation}
忽略因子二，涨落尺度的复数记法为
\[
x(t)\sim\sqrt{\frac{\hbar}{\omega}}\,e^{i\omega t}.
\]
求导则给出
\begin{equation}
\langle\dot x^{\,2}\rangle\sim\hbar\omega,
\label{eq:c10-zero-point-velocity}
\end{equation}
与平均动能和势能相等一致。冷却可以移除\(n>0\)的激发，但不能移除这些基态方差。

\section{阻尼、临界性与冻结}

将振子浸入黏性介质中。弹簧和质量块在蜂蜜中运动服从
\begin{equation}
\ddot x+\gamma\dot x+\omega^2x=0,
\label{eq:c10-damped}
\end{equation}
其中\(\gamma>0\)度量阻力。寻找\(x=e^{\alpha t}\)得到
\begin{equation}
\alpha^2+\gamma\alpha+\omega^2=0,
\end{equation}
因此
\begin{equation}
\alpha_\pm
=\frac{-\gamma\pm\sqrt{\gamma^2-4\omega^2}}{2}.
\label{eq:c10-damped-roots}
\end{equation}
这两个根提供了二阶微分方程所需的两组独立解。

\begin{classroomqa}
\audiencequestion{方程~\eqref{eq:c10-damped}中的阻力项是否应带负号？}
\lecturerresponse{阻力本身为\(F_{\rm drag}=-\gamma\dot x\)。牛顿方程为\(\ddot x=-\gamma\dot x-\omega^2x\)。将两个力项都移到左边，得到带正号的方程~\eqref{eq:c10-damped}。因此，微分方程中的符号与物理力的方向是一致的。\lecturetimestamp{00:59:38--01:00:10}}
\end{classroomqa}

如果\(\gamma^2<4\omega^2\)，则根为复数。定义
\[
\Omega=\sqrt{\omega^2-\frac{\gamma^2}{4}}.
\]
于是
\begin{equation}
x(t)=e^{-\gamma t/2}
\left(A\cos\Omega t+B\sin\Omega t\right).
\label{eq:c10-underdamped}
\end{equation}
振子仍会反复穿过平衡点，但其包络衰减。这就是欠阻尼 regime。

如果\(\gamma^2>4\omega^2\)，则方程~\eqref{eq:c10-damped-roots}的两个根均为实数且为负：
\begin{equation}
x(t)=A e^{\alpha_+t}+B e^{\alpha_-t}.
\label{eq:c10-overdamped}
\end{equation}
不存在持续的振荡；两个分量都衰减。这就是过阻尼区制。在
\begin{equation}
\gamma=2\omega
\label{eq:c10-critical}
\end{equation}
处两根重合，标志着临界阻尼，即两种行为之间的过渡。

一个极限情形尤为重要。没有回复力时，\(\omega=0\)，
\[
\alpha(\alpha+\gamma)=0,
\]
所以
\begin{equation}
x(t)=c_0+c_1e^{-\gamma t}.
\label{eq:c10-no-restoring}
\end{equation}
速度消亡，但坐标不必回到零。它冻结在由初始条件选定的常数\(c_0\)处。一块没有连接弹簧的石头被推着穿过蜂蜜，会在阻力耗尽其运动的地方停下来。

现在让频率缓慢变化。从\(\omega\gg\gamma\)开始，因此运动是振荡的。在保持阻尼尺度大致不变的同时减小\(\omega(t)\)。振荡减慢，越过\(\gamma=2\omega\)，变为过阻尼，并当\(\omega\to0\)时趋近无外力形式。等价地，也可以相对于回复力增大阻尼。在宇宙学应用中，决定性的变化将是有效频率的下降，而哈勃阻尼几乎保持不变。最终的冻结值不能仅由过渡点推断；它来自完整的早期演化和初始量子态。

\section{膨胀宇宙中的标量波}

回到暴胀子\(\phi\)。假定其势包含一个高而宽的平台，随后陡峭下降到一个较低的极小值。在平台上，重要的近似是
\begin{equation}
V'(\phi)\simeq0,
\label{eq:c10-flat-potential}
\end{equation}
而不是\(V(\phi)\simeq0\)。斜率很小，而能量密度仍可大到足以驱动暴胀。这种形状的具体起源并未给出；它是该模型的一个唯象要求。

\begin{figure}[htbp]
\centering
\includegraphics[width=0.88\textwidth]{lecture_10_figure_02.png}
\caption{左侧仍可见暴胀子平台和陡崖，而右侧开始写出平直空间波动方程。黑板上可靠地支持近似条件\(V'(\phi)=0\)以及一个二阶导数项的开头，但讲演者遮住了未完成的方程；方程~\eqref{eq:c10-flat-wave}是根据口头推导重建的。}
\label{fig:c10-potential-wave}
\lectureframe{01:17:17.865}
\end{figure}

对于平直时空中的近自由标量场，
\begin{equation}
\frac{\partial^2\phi}{\partial t^2}
-\nabla^2\phi=0.
\label{eq:c10-flat-wave}
\end{equation}
非平坦的势会在右边加上\(-V'(\phi)\)。对于平台近似，我们将其略去。如果场是完全均匀的，空间导数会消失。在膨胀宇宙中，剩下的均匀方程将是
\[
\ddot\phi+3H\dot\phi=0,
\]
即上一讲中的哈勃摩擦方程。

\begin{figure}[htbp]
\centering
\includegraphics[width=0.88\textwidth]{lecture_10_figure_03.png}
\caption{左侧标出了平台，而右侧出现了波动方程的相继阶段。上方的平直空间算符清晰可读；下方一行写到了\(+3H\)但未完成，因此完整的阻尼项是根据当时的解释补上的。}
\label{fig:c10-hubble-board}
\lectureframe{01:18:57.554}
\end{figure}

\begin{figure}[htbp]
\centering
\begin{coursefitfigure}
\begin{tikzpicture}[x=1cm,y=0.9cm,>=Latex]
  \draw[->] (-0.2,0) -- (9.3,0) node[right] {\(\phi\)};
  \draw[->] (0,-0.2) -- (0,4.6) node[above] {\(V(\phi)\)};
  \draw[very thick,blue!55!black,smooth]
    plot coordinates {(0.25,3.9) (1.4,3.88) (2.8,3.8) (4.2,3.6) (5.3,3.25) (6.0,2.6) (6.5,1.45) (6.9,0.45) (7.5,0.28) (8.5,0.7)};
  \draw[dashed] (6.45,0) -- (6.45,2.0);
  \draw[->,thick,red!65!black] (2.0,4.25) -- (3.1,3.82);
  \node[anchor=south west] at (0.8,4.0) {slow-roll plateau: \(V'\simeq0\)};
  \node[anchor=west] at (6.55,2.05) {end of inflation};
  \node[anchor=north] at (7.45,0.22) {reheating region};
\end{tikzpicture}
\end{coursefitfigure}
\caption{整个论证中所用势的清晰重建。浅斜率允许慢滚，而\(V\)的高值使\(H\)几乎保持恒定；陡峭下降结束暴胀并释放储存的能量。}
\label{fig:c10-potential-reconstruction}
\end{figure}

对于非均匀场，必须区分坐标距离和物理距离。选择无量纲共动标签\(\mathbf x\)，于是无穷小固有距离为
\[
d\boldsymbol\ell=a(t)\,d\mathbf x.
\]
在固定时刻，
\[
\nabla_{\rm proper}=\frac{1}{a}\nabla_{\mathbf x},
\qquad
\nabla_{\rm proper}^2=\frac{1}{a^2}\nabla_{\mathbf x}^2.
\]
因此标量方程为
\begin{equation}
\ddot\phi+3H\dot\phi-\frac{1}{a^2}\nabla^2\phi=0,
\label{eq:c10-frw-wave}
\end{equation}
其中\(\nabla\)对共动标签求导。如果改为选择无量纲标度因子，则共动坐标带有长度；物理不变。

\begin{classroomqa}
\audiencequestion{方程~\eqref{eq:c10-frw-wave}中的空间项应当为正吗？}
\lecturerresponse{不应当。平直空间波动算符已经具有相对符号\(\partial_t^2-\nabla^2\)。膨胀以因子\(a^{-2}\)改变空间导数，但不会反转该符号。作傅里叶代换后，两个负号合并，模式方程获得预期的正回复项\(+k^2\phi_k/a^2\)。\lecturetimestamp{01:22:18--01:23:06}}
\end{classroomqa}

\section{每个傅里叶模式都是一个阻尼振子}

一次研究一个空间波：
\begin{equation}
\phi(\mathbf x,t)
=\phi_{\mathbf k}(t)e^{i\mathbf k\cdot\mathbf x}.
\label{eq:c10-mode-ansatz}
\end{equation}
复指数只是正弦和余弦的记账方式。任意实轮廓通过叠加模式来构造，其中\(\mathbf k\)与\(-\mathbf k\)之间满足适当的实性关系。

时间导数只作用于\(\phi_{\mathbf k}(t)\)，而
\[
\nabla^2e^{i\mathbf k\cdot\mathbf x}
=-k^2e^{i\mathbf k\cdot\mathbf x}.
\]
代入方程~\eqref{eq:c10-frw-wave}得到
\begin{equation}
\ddot\phi_{\mathbf k}
+3H\dot\phi_{\mathbf k}
+\frac{k^2}{a^2}\phi_{\mathbf k}=0.
\label{eq:c10-mode-equation}
\end{equation}
因此每个波数都像一个阻尼振子，其
\begin{equation}
\gamma=3H,
\qquad
\omega_k(t)=\frac{k}{a(t)}.
\label{eq:c10-mode-identification}
\end{equation}
物理波数\(k/a\)随宇宙膨胀而红移。因而有效回复力\(k^2/a^2\)减小并趋于零。

一个起始于短物理波长的模式具有 \(\omega_k\gg H\) 并振荡。它的波长被拉伸，频率下降，哈勃阻尼最终占主导。随后它趋近于一个不振荡的值，而不是被拉回零。不同的 \(\mathbf k\) 模式在不同的时间发生跨越：在固定 \(a\) 下，较大的 \(k\) 对应较短的波长，必须等待更长时间的膨胀，其物理频率才会降至哈勃尺度。

\section{视界退出与冻结的真空涨落}

以振子的临界条件作为量级为一的估计，
\[
\gamma=2\omega_k,
\]
式~\eqref{eq:c10-mode-identification}给出
\begin{equation}
3H=\frac{2k}{a},
\qquad
a_{\rm cross}=\frac{2k}{3H}.
\label{eq:c10-critical-mode}
\end{equation}
数值因子并非概念要点。跨越发生在物理波数 \(k/a\) 为 \(H\) 量级之时。

令
\[
\lambda_{\rm com}=\frac{2\pi}{k},
\qquad
\lambda_{\rm phys}=a\lambda_{\rm com}
=\frac{2\pi a}{k}.
\]
量纲检验确定了 \(a\) 的位置：在无量纲的 \(\mathbf x\) 下，\(k\) 与 \(\mathbf k\cdot\mathbf x\) 均为无量纲，而 \(a\) 与 \(\lambda_{\rm phys}\) 带有长度量纲。因此式~\eqref{eq:c10-critical-mode}蕴含
\begin{equation}
\lambda_{\rm phys}\sim H^{-1},
\label{eq:c10-horizon-crossing}
\end{equation}
至多相差诸如 \(2\)、\(3\) 和 \(2\pi\) 之类的因子。

哈勃定律 \(v=Hd\) 将距离 \(H^{-1}\) 与量级为 \(c=1\) 的退行速度联系起来。在暴胀期间，\(H\) 很大且近乎恒定，因此这一哈勃长度是微观的，尽管它扮演着因果尺度的角色。将这一事件称为 \emph{视界退出} 是约定俗成的：没有任何物质客体穿过一堵物理的墙。该模式的物理波长增长到超过近似固定的哈勃尺度，其动力学从快速振荡变为几乎冻结的振幅。

\begin{figure}[htbp]
\centering
\begin{coursefitfigure}
\begin{tikzpicture}[x=1cm,y=0.85cm,>=Latex]
  \draw[->] (0,0) -- (9.2,0) node[right] {cosmic time};
  \draw[->] (0,0) -- (0,5.0) node[above] {physical length};
  \draw[dashed,thick] (0.25,2.75) -- (8.75,2.75)
    node[right] {\(H^{-1}\)};
  \draw[very thick,blue!55!black]
    (0.35,0.55) .. controls (3.3,0.75) and (5.3,1.55) .. (8.5,4.55);
  \node[blue!55!black,anchor=west] at (6.65,4.55) {\(\lambda_{\rm phys}=2\pi a/k\)};
  \draw[densely dotted] (6.05,0) -- (6.05,2.75);
  \fill (6.05,2.75) circle (2pt);
  \node[anchor=north] at (6.05,0) {horizon exit};
  \node[align=center] at (2.6,1.7) {oscillating\\sub-Hubble mode};
\end{tikzpicture}
\end{coursefitfigure}
\caption{在近乎指数膨胀期间，物理波长增长，而哈勃长度保持近似固定。它们的交叉标志着从快速振荡到冻结模式的转变，至多相差量级为一的因子。}
\label{fig:c10-horizon-reconstruction}
\end{figure}

真空涨落使暴胀无法简单地抹去每一个波。每一个短波长模式都从零点运动开始。膨胀将一个模式拉伸到哈勃尺度并将其冻结，而更短的真空模式则继续振荡。它们依次被拉伸并冻结。如此反复，便建立起一个场分布，其中包含跨越广泛波长和取向范围的冻结贡献。其结果是随机的，但并非无结构的：其统计性质由量子态和暴胀背景决定。

\begin{figure}[htbp]
\centering
\includegraphics[width=0.88\textwidth]{lecture_10_figure_04.png}
\caption{一幅横跨整个黑板的场波长叠加图，包含长波与短波，并在标签 \(\gamma\) 和 \(3H\) 旁附有一个较小的阻尼示意图。该图记录了定性的模式图像；它并非用于读出一个精确的函数。}
\label{fig:c10-mode-board}
\lectureframe{01:45:28.476}
\end{figure}

\begin{classroomqa}
\audiencequestion{模式是在一个点上冻结，还是在一个波长上冻结，而 \(H\) 本身是否在指数增长？}
\lecturerresponse{当波长达到哈勃尺度时它们冻结。在近乎平坦的平台区，\(V\) 因而 \(H\) 近似恒定；是指标因子 \(a(t)\)，从而物理波长 \(a/k\)，在指数增长。每一个共动模式在各自的时间达到条件 \(k/a\sim H\)。\lecturetimestamp{01:45:06--01:45:39}}
\end{classroomqa}

\begin{classroomqa}
\audiencequestion{这些是电磁波吗？}
\lecturerresponse{不是。它们是暴胀子标量场的涨落。电磁真空涨落提供了一个熟悉的类比，但式~\eqref{eq:c10-mode-equation}中的变量是 \(\phi\)。\lecturetimestamp{01:45:39--01:46:04}}
\end{classroomqa}

\begin{classroomqa}
\audiencequestion{是否存在一个连续的波长谱？}
\lecturerresponse{对于无限空间的理想化情形，是的。场由傅里叶积分表示，
\begin{equation}
\phi(\mathbf x,t)
=\int\frac{d^3k}{(2\pi)^3}\,
\phi_{\mathbf k}(t)e^{i\mathbf k\cdot\mathbf x}.
\label{eq:c10-fourier-integral}
\end{equation}
每一个波数都经历相同的红移和冻结机制。\lecturetimestamp{01:46:04--01:46:34}}
\end{classroomqa}

\begin{classroomqa}
\audiencequestion{有限宇宙是否只有离散的模式？}
\lecturerresponse{是的。有限空间体积上的边界条件将允许的波数量子化。对于远大于所讨论尺度的体积，间距可以如此之细，以至于求和可由式~\eqref{eq:c10-fourier-integral}中的积分很好地近似。连续图像是一种理想化，而非一种不同的局域机制。\lecturetimestamp{01:46:34--01:47:01}}
\end{classroomqa}

如今两种运动并存。空间平均场 \(\bar\phi(t)\) 缓慢而经典地沿势能向下滚动。在其之上叠加着一个小的量子涨落，
\[
\phi(\mathbf x,t)=\bar\phi(t)+\delta\phi(\mathbf x,t),
\]
它由冻结的模式组装而成。涨落能量本身太小，不足以成为所观测到的物质对比。其重要性在于它扰动了暴胀结束的时间。

完整的单场量子计算使这一振子论证更加精确。一个轻标量场在每个波数的对数区间内获得的典型涨落量级为
\begin{equation}
\delta\phi_{\rm exit}\sim\frac{H}{2\pi}.
\label{eq:c10-field-amplitude}
\end{equation}
要推导其归一化，需要在暴胀背景中对场模式进行量子化。振子图像已经提供了关键机制：真空涨落即使在基态中也存在，而膨胀将它们转化为一种看起来像经典的冻结图案。

\section{悬崖将场噪声转化为密度噪声}

首先抑制\(\delta\phi\)。于是在任意固定时刻，场在各处都只有一个值；当空间\(x\)向上、场值\(\phi\)向右画出时，这在黑板上表示为一条竖直线。随着时间推进，这条线向悬崖移动。整个宇宙一起到达下落处，暴胀子势能变成粒子、辐射和热，指数膨胀结束，由此产生的能量密度开始随普通膨胀而稀释。

\begin{figure}[htbp]
\centering
\includegraphics[width=0.88\textwidth]{lecture_10_figure_05.png}
\caption{黑板上的空间-场图。竖直轴是\(x\)，水平方向是\(\phi\)，淡淡的近竖直剖面标出相继的场位形，右侧带阴影的边界就是暴胀结束处的悬崖。}
\label{fig:c10-cliff-board}
\lectureframe{01:50:54.502}
\end{figure}

\begin{figure}[htbp]
\centering
\begin{coursefitfigure}
\begin{tikzpicture}[x=1cm,y=0.82cm,>=Latex]
  \draw[->] (0,0) -- (9.2,0) node[right] {\(\phi\)};
  \draw[->] (0,0) -- (0,5.2) node[above] {space \(x\)};
  \draw[thick,gray!60] (8.2,0.15) rectangle (8.55,4.85);
  \foreach \y in {0.35,0.85,...,4.55} {
    \draw[gray!60] (8.2,\y) -- (8.55,{\y+0.25});
  }
  \node[anchor=south,rotate=90] at (8.38,2.5) {cliff};
  \draw[thick,densely dashed,gray!70] (3.2,0.2) -- (3.2,4.75);
  \draw[thick,blue!45!black,smooth]
    plot coordinates {(5.0,0.2) (5.12,0.75) (4.9,1.25) (5.18,1.8) (4.94,2.35) (5.2,2.9) (4.88,3.45) (5.12,4.05) (5.0,4.75)};
  \draw[very thick,red!65!black,smooth]
    plot coordinates {(7.15,0.2) (7.32,0.7) (7.0,1.2) (7.42,1.75) (7.1,2.3) (7.5,2.85) (7.02,3.35) (7.35,4.0) (7.16,4.75)};
  \draw[->,thick] (3.55,2.5) -- (4.55,2.5);
  \draw[->,thick] (5.55,2.5) -- (6.65,2.5);
  \node[align=center,anchor=south,font=\small] at (3.2,4.8) {uniform\\profile};
  \node[align=center,anchor=south,font=\small] at (5.0,4.8) {frozen\\ripple};
  \node[align=center,anchor=south,font=\small] at (7.2,4.8) {near\\reheating};
\end{tikzpicture}
\end{coursefitfigure}
\caption{空间-场图像的清晰重建。平均场向悬崖移动，同时一个冻结的空间波纹持续存在，因此不同位置在不同时刻略微不同地到达暴胀结束。}
\label{fig:c10-cliff-reconstruction}
\end{figure}

恢复冻结的起伏。这条线现在略微呈波浪状：相邻位置的\(\phi\)值略有不同。设想一排人试图肩并肩地向悬崖行进。有几个人稍微靠前，有几个人稍微靠后。靠前的人先掉下去。

这种时间偏移才是密度反差的真正来源。先掉下去的区域先重加热。它的势能被倾倒入粒子和热中，随后在暴胀后时期宇宙膨胀时，这部分能量有时间稀释。附近稍晚掉下去的区域稍晚重加热，因此当两者相比较时，较晚的区域更密，仅仅是因为它稀释的时间更短。因此，第一组穿越留下了相对致密的区域。当这些区域膨胀时，较短波长的波纹又把另一组位置带过边缘，并在较大图案内部形成较小的区域。在冻结模式的整个层级上重复这一点，就产生了结构中的结构。因此，微小的场涨落之所以重要，并不是因为它自身的能量，而是因为它改变了局部重加热发生的时间。

这种时间关系可以示意性地写出。如果均匀场以速度\(\dot{\bar\phi}\)到达终值\(\phi_{\rm end}\)，那么一个小的场位移会使局部结束时间改变
\begin{equation}
\delta t_{\rm end}
\simeq-\frac{\delta\phi}{\dot{\bar\phi}}.
\label{eq:c10-time-delay}
\end{equation}
滚动越慢，同样的\(\delta\phi\)就对应越大的延迟。在这段额外时间内，较早的区域膨胀并稀释得更多。直到量级为一的转移因子，曲率或密度扰动因此由下式控制
\begin{equation}
\mathcal R
\sim H\,\delta t_{\rm end}
\sim-\frac{H\,\delta\phi}{\dot{\bar\phi}}.
\label{eq:c10-curvature-estimate}
\end{equation}
将其与方程~\eqref{eq:c10-field-amplitude}结合，就给出熟悉的慢滚尺度
\begin{equation}
\mathcal P_{\mathcal R}^{1/2}
\sim\frac{H^2}{2\pi|\dot{\bar\phi}|},
\label{eq:c10-slow-roll-amplitude}
\end{equation}
而精确的约定以及随后向\(\delta\rho/\rho\)的转移则需要更完整的扰动计算。场滚动得越慢，相邻悬崖穿越之间的时间就越长，其间发生的稀释也越大。因此，当\(\lvert\dot{\bar\phi}\rvert\)更小时，同样的场涨落会留下更大的密度印记。

许多模式在不同时间越过哈勃尺度，因此许多大小不同的区域在略微不同的时间重加热。由此产生的密度场包含在尺度上嵌套的过密和欠密区域。其统计性质取决于暴胀势的高度、斜率和曲率。合适的慢滚模型自然会产生接近
\[
\frac{\delta\rho}{\rho}\sim10^{-5},
\]
的近标度不变振幅，而详细的CMB谱使得少量宇宙学参数可以在许多角尺度上受到检验。

\section{从原初种子回到天空}

现在我们可以回到位于一组过去光壳中心的观测者。最后散射球面穿过三维原初密度场。该球面与不同引力势相交的方向会产生略微不同的微波温度。WMAP 展示了该球面上的二维图案。这张图既不是时间零点的照片，也不是球面内部所有结构的投影；它是微波光子最后散射的那个特定球面上的温度结构。

在球面内部，同样的原初扰动继续演化。引力放大过密区，类焦散流使片状结构和纤维状结构变得尖锐，气体聚集，星系形成。因此，CMB 图案和后来的宇宙网是同一历史的两个阶段：冻结的暴胀涨落首先表现为微小的密度种子，后来成为可见结构。

\begin{classroomqa}
\audiencequestion{我们今天看到的微波背景代表的是哪个时期？}
\lecturerresponse{这些光子已经旅行了大约 \(13\) 亿年，但它们是在热大爆炸之后几十万年被释放的，当时电子和质子结合成中性氢，宇宙变得透明。暴胀结束得要早得多。重加热产生了热物质和辐射；宇宙随后膨胀并冷却，直到最后散射。对于这一论证，稳健的陈述是数量级---几十万年---，而不是精确的复合日期。 \lecturetimestamp{02:00:38--02:02:19}}
\end{classroomqa}

\section{一片天空与宇宙方差}

\begin{classroomqa}
\audiencequestion{暴胀是否预言了 WMAP 中报道的优选方向或北--南不对称？}
\lecturerresponse{从本文发展的简单暴胀图景中，并不会得出这样的优选方向。所报道的大角度异常的地位当时并不确定，前景和统计涨落是现实关切，负责任的评估需要专业文献，而不是耸人听闻的概述。暴胀预言的是一个系综的统计性质；判断一个看起来不寻常的天空特征是否是新物理则更难。 \lecturetimestamp{02:02:19--02:03:26}}
\end{classroomqa}

小角度特征在天空中会出现许多次。它们提供了许多样本，因此可以对其分布和各向同性进行统计检验。四极矩尺度的特征则不同：在可观测球面上只能容纳少数几个独立的大角度模式。因此，即使基础理论正确，它们测得的功率也可能与系综平均值相差很大。这种不可约的采样不确定性就是\emph{宇宙方差}。

\begin{classroomqa}
\audiencequestion{那么几个低多极矩的表观对齐呢？}
\lecturerresponse{它可能是一条线索、一个前景或分析伪影，或者一个不太可能但被允许的实现。困难在于，我们无法从同一统计系综中再抽取另一个可观测宇宙。一枚硬币在六次抛掷中给出四次正面，并不一定有偏；重复多批抛掷才能判定，但宇宙学在最大尺度上只提供了一批。这个异常也许值得研究，然而仅凭一片天空本身无法提供决定性显著性检验所需的重复试验。 \lecturetimestamp{02:05:31--02:07:36}}
\end{classroomqa}

\section{从量子运动到宇宙结构}

三个普适常数定义了普朗克尺度，并揭示出熟悉的粒子物理距离量子引力有多远。微波天空提供了一个精确的宇宙学目标：接近各向同性、涨落幅度约为 \(10^{-5}\)，以及结构分布在许多尺度上。量子力学确保每个场模式都有零点涨落。暴胀拉伸每个暴胀子模式，直到其物理波长达到 \(H^{-1}\) 的量级，此时哈勃阻尼抑制其振荡，并留下一个有效冻结的幅度。

均匀的暴胀子继续在那个冻结的噪声之下滚动。当该场到达其平台末端时，相邻区域不会在完全相同的时间重加热。较早的区域稀释得更久；较晚的区域则保持更密。这将 \(\delta\phi\) 转化为原初密度和曲率图案。随后引力将该图案放大为星系和宇宙网，而最后散射球面则在微波天空中保存了它较早的状态。因此，微观零点运动、宇宙膨胀和交错重加热并不是彼此分离的故事。它们共同解释了为什么一个暴胀宇宙既能极其光滑，又具有足以让我们存在的结构。
